\documentclass[12pt,letterpaper]{article}

\usepackage[margin=1in]{geometry}
\usepackage{amsmath,amssymb}
\usepackage{fancyhdr}

\setlength{\parindent}{0pt}
\setlength{\parskip}{0.8em}
\setlength{\headheight}{15pt}

\pagestyle{fancy}
\fancyhf{}
\fancyhead[L]{FIT3155 SAMPLE EXAM PAPER}
\fancyfoot[C]{Page \thepage{} of 22}
\renewcommand{\headrulewidth}{0pt}
\renewcommand{\footrulewidth}{0pt}

\newcommand{\marksbox}[1]{%
  \vfill
  \begin{flushright}
    \fbox{\parbox[c][1.1cm][c]{1.3cm}{\centering #1}}
  \end{flushright}
}

\newcommand{\questiontitle}[2]{%
  \textbf{Question #1}\hfill\textbf{(#2 marks)}
}

\newcommand{\workingpage}{%
  \newpage
  \begin{center}
    (Page intentionally left blank for working space)
  \end{center}
}

\begin{document}

\begin{titlepage}
  \thispagestyle{empty}
  \vspace*{0.18\textheight}
  {\Large FIT3155 \hfill SAMPLE EXAM PAPER\par}
  \vspace{2.2cm}
  \begin{center}
    {\LARGE\bfseries FIT3155 SAMPLE EXAM PAPER\par}
    \vspace{0.8cm}
    {\large Generated Sample Exam 02\par}
  \end{center}
\end{titlepage}

\setcounter{page}{2}

\section*{INSTRUCTIONS}

\begin{itemize}
  \item This exam is worth 60 marks.
\end{itemize}

\begin{center}
  \textbf{General exam technique}
\end{center}

\begin{itemize}
  \item Some students lose marks by not attempting all questions.

  \item Suppose you are likely to get 3/6 on a question after spending say 10 minutes on it.
  Spending another 10 minutes on the same question may gain you at most 3 additional
  marks. In contrast, spending those 10 minutes on a new question could earn you another
  6 marks. The key point is: do not get stuck on one question.

  \item A good strategy is to first answer the questions you are most confident about before
  attempting the others. In other words, answering questions strictly serially during an
  exam can be suboptimal.

  \item Do not waste time answering an unrelated question in the hope of gaining partial marks.
  Marking will be based on the question that was actually asked.

  \item Where necessary, justify your answers with clear and concise explanations, without being
  overly wordy.
\end{itemize}

\newpage

\questiontitle{1}{6}

A \emph{border} of a string $S[1 \ldots n]$ is any \emph{proper} prefix of $S$
(i.e., a prefix of length strictly less than $n$) that is also a suffix of $S$.

Example: $S[1 \ldots 7] = \texttt{abacaba}$ has exactly two borders,
\texttt{a} (of length 1) and \texttt{aba} (of length 3).

Write pseudocode describing an efficient algorithm that accepts any input
string $S[1 \ldots n]$ and outputs the lengths of \emph{all} borders of $S$.
State the worst-case time complexity of your algorithm.

(For this problem, you may assume that \texttt{zalg} is available as a built-in
function that takes a string as input and returns the array of its Z-values.
Do not write pseudocode for \texttt{zalg}.)

\marksbox{6}

\workingpage
\newpage

\questiontitle{2}{6}

In the Boyer--Moore algorithm for exact matching of the pattern
$pat[1 \ldots m]$ against the text $txt[1 \ldots n]$, consider a general
iteration in which the region $txt[j \ldots j+m-1]$ of the text is being
compared with $pat$, and a right-to-left scan results in a mismatch at
position $1 \leq k \leq m$ of the pattern, against some character $x$ in the
text.

Let $R(x)$ denote the position of the \emph{rightmost} occurrence of the
character $x$ in $pat$, and suppose that $R(x) < k$. Under the bad character
rule, $pat$ is shifted $k - R(x)$ positions to the right, so that the
occurrence of $x$ at $pat[R(x)]$ becomes aligned with the mismatched character
$x$ in the text.

Prove that this bad character shift is safe, i.e., such a shift can never skip
$pat$ past any valid occurrence in $txt$.

\marksbox{6}

\workingpage
\newpage

\questiontitle{3}{6=3+3}

Assume throughout that $\$$ is a unique terminal symbol that is
lexicographically smaller than every other character.

\begin{enumerate}
  \item[(i)] Compute the Burrows--Wheeler Transform of the string
  \[
    S = \texttt{cocoa\$}.
  \]
  Show the sorted rotations (or sorted suffixes) you used to obtain your
  answer.

  \item[(ii)] The Burrows--Wheeler Transform of some unknown string $T\$$ is
  \[
    L[1 \ldots 6] = \texttt{abbaa\$}.
  \]
  Recover $T$. Show your working, e.g., the last-to-first (LF) column mapping
  you used at each step of the inversion.
\end{enumerate}

\marksbox{6}

\workingpage
\newpage

\questiontitle{4}{6=3+3}

These questions relate to Ukkonen's linear-time suffix tree construction
algorithm for a string $S[1 \ldots n]$.

\begin{enumerate}
  \item[(i)] Suppose suffix extension $j$ of some phase $i$ creates a new leaf
  numbered $j$ (using Rule 2). Prove that in \emph{every} subsequent phase
  $i' > i$, suffix extension $j$ is completed by Rule 1 (the `leaf-edge
  extension rule'), i.e., ``once a leaf, always a leaf''.

  \item[(ii)] In the \emph{rapid leaf extension} trick, the label of every
  leaf edge is stored as a pair
  \[
    (start,\ \texttt{globalEnd}),
  \]
  where \texttt{globalEnd} is a single mutable integer shared by all leaf
  edges. Describe how incrementing \texttt{globalEnd} implements all Rule 1
  extensions of a phase in $O(1)$ total time, and explain why your result
  from part (i) guarantees that this trick is sound.
\end{enumerate}

\marksbox{6}

\workingpage
\newpage

\questiontitle{5}{6=3+3}

\begin{enumerate}
  \item[(i)] Decode the following Elias-Omega codeword to recover the positive
  integer $N$ that it encodes:
  \[
    \texttt{0011101}
  \]
  Show each component you read, and how it determines the length of the next
  component.

  \item[(ii)] Encode the following string into
  $\langle \textit{offset}, \textit{length}, \textit{next char} \rangle$
  triples using the LZ77 algorithm:
  \[
    \texttt{abcabcabcd}
  \]
  Assume that the search window size and the lookahead buffer size are both
  15. List every triple your encoder emits, in order.
\end{enumerate}

\marksbox{6}

\workingpage
\newpage

\questiontitle{6}{6=3+3}

Let $u$ and $v$ be two $n$-bit positive integers, where $n$ is a power of two.
Write
\[
  u = 2^{n/2}\,U_1 + U_0, \qquad v = 2^{n/2}\,V_1 + V_0,
\]
where $U_1, U_0, V_1, V_0$ are $n/2$-bit integers.

\begin{enumerate}
  \item[(i)] Show how the product $u \cdot v$ can be computed using only
  \emph{three} multiplications of $n/2$-bit integers, together with additions,
  subtractions, and bit shifts (Karatsuba's method). State the identity you
  use explicitly.

  \item[(ii)] Write down the recurrence describing the running time $T(n)$ of
  the resulting divide-and-conquer multiplication algorithm, and solve it to
  obtain a tight asymptotic bound on $T(n)$.
\end{enumerate}

\marksbox{6}

\workingpage
\newpage

\questiontitle{7}{6=3+3}

A Fibonacci heap $H$ consists of three min-heap-ordered trees, whose roots
appear in the root list in the following order (left to right):

\begin{itemize}
  \item a tree rooted at $4$, where node $4$ has children $9$ and $11$, and
  node $9$ has a single child $12$;
  \item a tree rooted at $6$, where node $6$ has a single child $8$;
  \item a single node $18$.
\end{itemize}

Node $9$ is \emph{marked}; no other node is marked. $H.min$ points to node
$4$.

Use the following conventions: nodes cut from a tree are appended to the right
end of the root list in the order they are cut; during consolidation the root
list is processed left to right; and when two trees of equal degree are
merged, the root with the smaller key becomes the parent.

\begin{enumerate}
  \item[(i)] Perform \textsc{Decrease-Key}$(12 \rightarrow 3)$ on $H$. Draw
  the resulting heap, indicating $H.min$ and any marked nodes.

  \item[(ii)] On the heap resulting from part (i), perform
  \textsc{Extract-Min}. Draw the resulting heap after consolidation,
  indicating $H.min$ and any marked nodes.
\end{enumerate}

\marksbox{6}

\workingpage
\newpage

\questiontitle{8}{6}

A $k$-bit binary counter data structure supports an increment by 1 operation
(modulo $2^k$), starting with all $k$ bits of the counter initialized to 0.

Using the \emph{potential} method of amortized analysis, determine the
amortized time complexity of an increment operation for this data structure.

State your potential function $\Phi$ explicitly, verify that it satisfies the
required boundary conditions, and show the calculation of the amortized cost
of a general increment.

\marksbox{6}

\workingpage
\newpage

\questiontitle{9}{6=3+3}

Insert the keys
\[
  10,\ 20,\ 30,\ 40,\ 50,\ 25,\ 35,\ 5
\]
in this order into an initially empty B-tree with minimum degree $t = 2$
(i.e., every node holds at most $2t - 1 = 3$ keys).

Use \emph{proactive} (top-down) splitting: while travelling from the root
towards the insertion leaf, any full node encountered is split before it is
stepped into.

\begin{enumerate}
  \item[(i)] Draw the B-tree immediately after the insertion of key $40$.

  \item[(ii)] Draw the final B-tree after all eight insertions. Clearly
  indicate the keys in every node.
\end{enumerate}

\marksbox{6}

\workingpage
\newpage

\questiontitle{10}{6}

Solve the following optimization problem using the tableau simplex algorithm.
Show the tableau after each pivot, and state the optimal values of $x$, $y$,
and $z$.

\[
\begin{array}{rl}
\text{Maximize} & z = 2x + 3y \\
\text{Subject to} & x + y \leq 4, \\
                  & x + 3y \leq 6, \\
                  & x \geq 0,\quad y \geq 0.
\end{array}
\]

\marksbox{6}

\workingpage

\vfill
\begin{center}
  End of Exam
\end{center}

\end{document}
